\section{De las redes totalmente conectadas a la convolución: el aprovechamiento de la estructura espacial}

\subsection{El problema de procesar imágenes con redes totalmente conectadas}

En la primera clase aprendimos las redes neuronales totalmente conectadas (Fully Connected Neural Network). Si se usa directamente una red totalmente conectada para procesar una imagen, hay que \textbf{aplanar} (flatten) primero la imagen bidimensional en un vector unidimensional:

\begin{figure}[H]
\centering
\includegraphics[width=0.9\textwidth]{slides-images/slide-22.jpg}
\caption{El problema de procesar imágenes con redes totalmente conectadas: la operación de aplanado destruye toda la información espacial y produce una cantidad enorme de parámetros.\protect\footnotemark}
\end{figure}
\footnotetext{Diapositivas, página 22.}

\begin{warningbox}{Los dos grandes defectos de las redes totalmente conectadas aplicadas a imágenes}
\begin{enumerate}
    \item \textbf{Pérdida total de la información espacial}: al aplanar la imagen bidimensional en un vector unidimensional, la relación espacial entre los píxeles queda completamente destruida. Píxeles que eran vecinos pueden quedar muy alejados entre sí después del aplanado.
    \item \textbf{Explosión del número de parámetros}: cada píxel se conecta con cada neurona de la capa oculta. Por ejemplo, una imagen en escala de grises de $256 \times 256$ tiene 65,536 píxeles; si la capa oculta tiene 1,024 neuronas, solo la primera capa requiere $65,536 \times 1,024 \approx 67$ millones de parámetros.
\end{enumerate}
\end{warningbox}

La pregunta clave es: \textbf{¿cómo aprovechar la estructura espacial de la entrada para orientar el diseño de la arquitectura de la red?}

\subsection{La idea de la conectividad local}

La idea central de la solución es sencilla: en lugar de que cada neurona se conecte a todos los píxeles de la imagen completa, se conecta únicamente a una \textbf{región local} (local patch) de la imagen de entrada.

\begin{figure}[H]
\centering
\includegraphics[width=0.85\textwidth]{slides-images/slide-23.jpg}
\caption{Aprovechamiento de la estructura espacial: se conecta una región local (patch) de la imagen de entrada a una sola neurona de la capa oculta. Cada neurona solo "ve" una pequeña parte de la imagen.\protect\footnotemark}
\end{figure}
\footnotetext{Diapositivas, página 23.}

Después, mediante una \textbf{ventana deslizante} (sliding window), la misma matriz de pesos se aplica repetidamente en distintas posiciones de la imagen. Esta es la idea central de la operación de \textbf{convolución} (Convolution).

\begin{figure}[H]
\centering
\includegraphics[width=0.85\textwidth]{slides-images/slide-24.jpg}
\caption{Ventana deslizante: se conecta una región local de la entrada a una neurona de la capa oculta, y la ventana deslizante define las conexiones. Pregunta clave: ¿cómo establecer los pesos para detectar una característica específica?\protect\footnotemark}
\end{figure}
\footnotetext{Diapositivas, página 24.}

\begin{importantbox}{Las tres propiedades centrales de la operación de convolución}
\begin{enumerate}
    \item \textbf{Extracción de características locales}: se usa un conjunto de pesos (filtro/filter) para extraer \textbf{características locales}
    \item \textbf{Múltiples filtros}: se usan \textbf{varios filtros} para extraer distintos tipos de características
    \item \textbf{Compartición de parámetros}: los parámetros de cada filtro se \textbf{comparten} espacialmente (spatial sharing), es decir, el mismo filtro se aplica en todas las posiciones de la imagen
\end{enumerate}
\end{importantbox}

\subsection{Resumen de la sección}

Al aplanar la imagen, la red totalmente conectada pierde la información espacial y produce una explosión del número de parámetros. La idea central de la convolución es la conectividad local junto con la compartición de parámetros, lo cual conserva la estructura espacial y reduce drásticamente el número de parámetros.

\newpage
